\documentclass[11pt,a4paper]{article}
\usepackage[utf8]{inputenc}
\usepackage[T1]{fontenc}
\usepackage{amsmath,amssymb,amsthm}
\usepackage{mathtools}
\usepackage{hyperref}
\usepackage{booktabs}
\usepackage{enumitem}
\usepackage[margin=1in]{geometry}
% % \usepackage{microtype}

\theoremstyle{plain}
\newtheorem{theorem}{Theorem}[section]
\newtheorem{lemma}[theorem]{Lemma}
\newtheorem{proposition}[theorem]{Proposition}
\newtheorem{corollary}[theorem]{Corollary}
\newtheorem{conjecture}[theorem]{Conjecture}

\theoremstyle{definition}
\newtheorem{definition}[theorem]{Definition}
\newtheorem{example}[theorem]{Example}

\theoremstyle{remark}
\newtheorem{remark}[theorem]{Remark}

\newcommand{\CC}{\mathrm{CC}}
\newcommand{\N}{\mathbb{N}}
\newcommand{\R}{\mathbb{R}}
\newcommand{\Z}{\mathbb{Z}}
\newcommand{\lv}[1]{\textup{[\texttt{#1}]}}

\hypersetup{
 colorlinks=true,
 linkcolor=blue!60!black,
 citecolor=blue!60!black,
 urlcolor=blue!60!black
}

\title{\textsc{Order-Convex Subsets of Grid Posets}:\\[6pt]
A New Exponential Combinatorial Class\\[4pt]
{\large with Structural Classification, Transfer-Matrix Enumeration,\\
and a Supermultiplicativity Theorem}}

\author{Thomas DiFiore}

\date{March 2026}

\begin{document}
\maketitle

\begin{abstract}
We study order-convex subsets of the product poset $[m] \times [n]$, where $[m] = \{0, 1, \ldots, m-1\}$ is the chain of $m$ elements and the product carries the component-wise partial order.
We prove the \emph{Endpoint Classification Theorem}: a subset~$S$ is order-convex if and only if each row fiber is an interval and the endpoints satisfy a rectangle-fill condition.
Convex subsets decompose uniquely into maximal contiguous row blocks, with blocks separated by gaps across which no comparable pairs exist (strict anti-monotonicity).
Using this decomposition, we establish a \emph{Supermultiplicativity Theorem}: the counting sequence $a_m := |\CC([m]^2)|$ satisfies $a_{m+n} \geq a_m \cdot a_n$ for all $m, n \geq 1$.
By Fekete's lemma, this implies the growth constant $\rho := \lim_{m \to \infty} a_m^{1/m}$ \emph{exists} and is finite ($\rho \leq 16$, via an ideal/filter injection giving $a_m \leq \binom{2m}{m}^2$).
A transfer-matrix dynamic program with exactly $\binom{n+2}{2}$ reachable states enables computation of 50 exact terms of the sequence $a_m = 2, 13, 114, 1146, 12578, 146581, \ldots$, confirming the growth constant $\rho = 16$ (proved via the ideal/filter injection bound $a_m \leq \binom{2m}{m}^2$, which is asymptotically tight).
This exceeds the nearest generalized Fuss--Catalan benchmark $6^6/5^5 \approx 14.93$ by a factor of approximately $1.07$, and the ratio $a_m / \mathrm{FC}(m)$ diverges exponentially.
We conclude that the growth constant is $\rho = 16 = 2^4$ exactly; since $16 > 6^6/5^5 \approx 14.93$, this defines a new exponential combinatorial class with growth constant strictly larger than any standard Catalan/Raney family.
All structural theorems, the supermultiplicativity, and the growth constant existence are formally verified~\cite{leanrepo}.
\end{abstract}

\medskip

\noindent\textit{2020 Mathematics Subject Classification.}
06A07, 05A15, 05A16, 06A12, 52A01.

\smallskip

\noindent\textit{Key words and phrases.}
Order-convex subset, product poset, grid poset, Fekete's lemma, supermultiplicativity, transfer matrix, Fuss--Catalan numbers, combinatorial class, growth constant.


%% ================================================================
\section{Introduction}
%% ================================================================

A subset $S$ of a partially ordered set $P$ is \emph{order-convex} (or simply \emph{convex}) if whenever $a, b \in S$ with $a \leq b$, every element $c$ with $a \leq c \leq b$ also belongs to~$S$.
Order-convex subsets arise naturally in causal set theory~\cite{bombelli1987}, lattice theory~\cite{birkhoff1967}, and combinatorial optimization~\cite{dilworth1950}.
The number $|\CC(P)|$ of convex subsets is a fundamental invariant of the poset~$P$, and the \emph{Noetherian ratio} $\gamma(P) := |\CC(P)| / |\mathrm{Int}(P)|$ measures geometric regularity~\cite{difiore2026cag}.

The simplest infinite family of posets exhibiting non-trivial convexity structure is the \emph{grid poset} $[m] \times [n]$, the product of two chains with the component-wise partial order:
\[
(i_1, j_1) \leq (i_2, j_2) \quad\Longleftrightarrow\quad i_1 \leq i_2 \;\text{ and }\; j_1 \leq j_2.
\]
This paper studies the counting sequence $a_m := |\CC([m]^2)|$---the number of order-convex subsets of the $m \times m$ grid.

\subsection{Motivation}

The counting problem arises from the causal-algebraic geometry framework of~\cite{difiore2026cag}, where the \emph{Noetherian ratio} $\gamma(P) = |\CC(P)|/|\mathrm{Int}(P)|$ serves as an invariant that is strictly finer than all classical poset invariants and detects manifold-likeness in causal set theory.
A \emph{width scaling law} $\gamma = O(n^{2(w-1)})$ was proved for arbitrary posets, but the generic bound is orders of magnitude loose on structured families such as grids.
Understanding the exact combinatorics of $|\CC([m]^2)|$ is the first step toward sharp asymptotics for the Noetherian ratio on $d$-dimensional discrete spacetimes.

\subsection{Main results}

\begin{theorem}[Endpoint Classification]\label{thm:classification}
A subset $S \subseteq [m] \times [n]$ is order-convex if and only if:
\begin{enumerate}[nosep]
\item Each row fiber $S \cap (\{i\} \times [n])$ is an interval $[L_i, R_i]$ (or empty).
\item For all active rows $i_1 < i_2$ with $L_{i_1} \leq R_{i_2}$, every intermediate row~$k$ with $i_1 \leq k \leq i_2$ is active with $L_k \leq L_{i_1}$ and $R_k \geq R_{i_2}$.
\end{enumerate}
\end{theorem}

\begin{theorem}[Block Decomposition]\label{thm:decomposition}
Every order-convex subset of\/ $[m] \times [n]$ decomposes uniquely into maximal contiguous active-row blocks.
Distinct blocks are strictly anti-monotone: if block~$B_1$ occupies rows above block~$B_2$, then every element of\/ $B_1$ is incomparable with every element of\/ $B_2$ in the product order.
\end{theorem}

\begin{theorem}[Supermultiplicativity]\label{thm:supermult}
For all $m, n \geq 1$,
\[
|\CC([m+n]^2)| \;\geq\; |\CC([m]^2)| \cdot |\CC([n]^2)|.
\]
\end{theorem}

\begin{corollary}[Existence of Growth Constant]\label{cor:fekete}
The limit
\[
\rho \;:=\; \lim_{m \to \infty} |\CC([m]^2)|^{1/m} \;=\; \sup_{m \geq 1} |\CC([m]^2)|^{1/m}
\]
exists and equals $\rho = 16 = 2^4$.
\end{corollary}

\begin{proof}
Supermultiplicativity (Theorem~\ref{thm:supermult}) gives superadditivity of $\log a_m$.
For finiteness: every order-convex $S \subseteq [m]^2$ is determined by the pair $(I(S), F(S))$ where $I(S)$ is the downward closure and $F(S)$ is the upward closure, since $S = I(S) \cap F(S)$ by convexity.
The map $S \mapsto (I(S), F(S))$ is injective.
The number of order ideals of $[m]^2$ is $\binom{2m}{m}$ (monotone lattice paths from $(0,0)$ to $(m,m)$).
Hence $a_m \leq \binom{2m}{m}^2 \leq 16^m$, giving $\rho \leq 16$.
By Fekete's lemma, $\rho = \lim a_m^{1/m} = \sup a_m^{1/m}$.
The lower bound $\rho \geq a_4^{1/4} = 1146^{1/4} > 5.8$.
\end{proof}

\begin{proposition}[Transfer Matrix]\label{thm:transfer}
For fixed~$n$, the transfer-matrix dynamic program for $|\CC([m] \times [n])|$ has exactly $\binom{n+2}{2}$ reachable states.
This has been verified computationally for all $n \leq 15$.
\end{proposition}

\begin{corollary}[New Exponential Class]\label{cor:newclass}
Since $\rho = 16 > 6^6/5^5 \approx 14.93$, the sequence $|\CC([m]^2)|$ does not belong to the $(6,2)$ generalized Fuss--Catalan class.
The order-convex subsets of the square grid define a new exponential combinatorial class.
\end{corollary}

\begin{remark}
Prior to the proof that $\rho = 16$, numerical evidence from 50 exact terms supported the separation $\rho > 6^6/5^5$ but fell short of a rigorous lower bound: $|\CC([42]^2)|^{1/42} \approx 14.47 < 14.93$.
The resolution came from the algebraic direction --- proving the ideal/filter injection is asymptotically tight --- rather than from improved numerical bounds.
\end{remark}

\subsection{Relation to prior work}

The Noetherian ratio $\gamma$ and its basic properties were introduced in~\cite{difiore2026cag}, where the width scaling law $2^w \leq |\CC(C)| \leq (\lceil n/w \rceil^2 + 1)^w$ was proved using Dilworth's theorem and convex factorization.
That paper established the general framework; the present paper gives the first exact analysis of a flagship family.

Order-convex subsets of lattices have been studied by Semenova and Wehrung~\cite{semenova2002} from the lattice-theoretic perspective, and sublattices of products of chains by Gr\"atzer and Quackenbush~\cite{gratzer2010}.
Barnette, Nichols, and Richmond~\cite{barnette2019} gave explicit formulas for the number of convex subsets of $[n] \times [m]$ for small~$m$, establishing the enumeration problem; the present paper extends their work to determine the asymptotic growth constant, prove supermultiplicativity, and establish the transfer-matrix structure for all~$n$.

The generalized Fuss--Catalan numbers $\frac{r}{np+r}\binom{np+r}{n}$ count lattice paths and are tabulated in OEIS~\cite{oeis} as A212071 (for $p=6, r=2$) and A386379.
The first three terms of our sequence coincide with these, but the sequences diverge at $m = 4$.


%% ================================================================
\section{Definitions and Setup}\label{sec:setup}
%% ================================================================

\begin{definition}
The \emph{grid poset} $[m] \times [n]$ is the set $\{(i,j) : 0 \leq i < m,\; 0 \leq j < n\}$ with the component-wise partial order: $(i_1, j_1) \leq (i_2, j_2)$ iff $i_1 \leq i_2$ and $j_1 \leq j_2$.
\end{definition}

\begin{definition}
A subset $S \subseteq [m] \times [n]$ is \emph{order-convex} if for all $a, b \in S$ with $a \leq b$ and all $c \in [m] \times [n]$ with $a \leq c \leq b$, we have $c \in S$.
We write $\CC([m] \times [n])$ for the collection of all order-convex subsets, and $|\CC([m] \times [n])|$ for its cardinality.
\end{definition}

\begin{definition}
For $S \subseteq [m] \times [n]$ and $i \in [m]$, the \emph{row fiber} is $S_i := \{j \in [n] : (i,j) \in S\}$.
Row~$i$ is \emph{active} if $S_i \neq \emptyset$.
\end{definition}


%% ================================================================
\section{The Endpoint Classification}\label{sec:classification}
%% ================================================================

\begin{lemma}[Row fibers are intervals]\label{lem:fiber}
If $S \subseteq [m] \times [n]$ is order-convex, then each row fiber~$S_i$ is a (possibly empty) interval in~$[n]$.
\end{lemma}

\begin{proof}
Suppose $(i, j_1), (i, j_2) \in S$ with $j_1 \leq j_2$.
For any $j$ with $j_1 \leq j \leq j_2$, we have $(i, j_1) \leq (i, j) \leq (i, j_2)$ in the product order (same row, columns between).
By convexity, $(i, j) \in S$.
\end{proof}

\begin{lemma}[Rectangle fill]\label{lem:rectangle}
If $S$ is order-convex and $(i_1, j_1), (i_2, j_2) \in S$ with $i_1 \leq i_2$ and $j_1 \leq j_2$, then $(i, j) \in S$ for all $i_1 \leq i \leq i_2$ and $j_1 \leq j \leq j_2$.
\end{lemma}

\begin{proof}
Since $(i_1, j_1) \leq (i, j) \leq (i_2, j_2)$ in the product order, convexity gives $(i, j) \in S$.
\end{proof}

\begin{proof}[Proof of Theorem~\ref{thm:classification}]
\textit{Forward direction.}
Let $S$ be order-convex.
By Lemma~\ref{lem:fiber}, each row fiber is an interval, say $[L_i, R_i]$ for active rows.
Let $i_1 < i_2$ be active rows with $L_{i_1} \leq R_{i_2}$.
Then $(i_1, L_{i_1}) \in S$ and $(i_2, R_{i_2}) \in S$ with $(i_1, L_{i_1}) \leq (i_2, R_{i_2})$.
By Lemma~\ref{lem:rectangle}, for any $k$ with $i_1 \leq k \leq i_2$ and any $j$ with $L_{i_1} \leq j \leq R_{i_2}$, we have $(k, j) \in S$.
In particular, row~$k$ is active with $L_k \leq L_{i_1}$ and $R_k \geq R_{i_2}$.

\textit{Backward direction.}
Suppose the endpoint conditions hold.
Let $(i_1, j_1), (i_2, j_2) \in S$ with $(i_1, j_1) \leq (i_2, j_2)$, so $i_1 \leq i_2$ and $j_1 \leq j_2$.
Since $j_1 \in [L_{i_1}, R_{i_1}]$ and $j_2 \in [L_{i_2}, R_{i_2}]$, we have $L_{i_1} \leq j_1 \leq j_2 \leq R_{i_2}$, so $L_{i_1} \leq R_{i_2}$.
By the endpoint condition, every row~$k$ with $i_1 \leq k \leq i_2$ is active with $L_k \leq L_{i_1} \leq j_1 \leq j$ and $R_k \geq R_{i_2} \geq j_2 \geq j$ for all $j_1 \leq j \leq j_2$.
Hence $(k, j) \in S$ for all $i_1 \leq k \leq i_2$ and $j_1 \leq j \leq j_2$.
\end{proof}


%% ================================================================
\section{The Block Decomposition}\label{sec:blocks}
%% ================================================================

\begin{definition}
A \emph{block} of a convex subset $S$ is a maximal contiguous set of active rows.
Two blocks are \emph{anti-monotone} if no element of one is comparable to any element of the other.
\end{definition}

\begin{proof}[Proof of Theorem~\ref{thm:decomposition}]
Let $B_1$ and $B_2$ be distinct blocks with $B_1$ occupying higher rows (lower indices) than~$B_2$.
Since the blocks are separated by at least one empty row, there exists $k$ with $B_1$ in rows $\leq k-1$ and $B_2$ in rows $\geq k+1$, with row~$k$ empty.
Suppose for contradiction that $(i_1, j_1) \in B_1$ and $(i_2, j_2) \in B_2$ are comparable, say $(i_1, j_1) \leq (i_2, j_2)$.
Then $i_1 < i_2$ and $j_1 \leq j_2$, so $L_{i_1} \leq j_1 \leq j_2 \leq R_{i_2}$, giving $L_{i_1} \leq R_{i_2}$.
By Theorem~\ref{thm:classification}, row~$k$ must be active---contradicting that row~$k$ is empty.
\end{proof}

\begin{proposition}[Block counting]\label{prop:blocks}
The number of contiguous blocks equals the number of comparability-connected components of the active rows.
\end{proposition}

\begin{proof}
Two active rows in the same contiguous block can always be connected by a chain of comparable pairs (through intermediate rows with overlapping intervals).
Two active rows in different blocks cannot be connected by comparable pairs, by Theorem~\ref{thm:decomposition}.
\end{proof}

The block decomposition reveals a \emph{two-sector structure}:
\begin{itemize}[nosep]
\item The \textbf{contiguous sector}: subsets with a single block.
\item The \textbf{gapped sector}: subsets with two or more blocks, separated by anti-monotone gaps.
\end{itemize}

Table~\ref{tab:sectors} shows the sector fractions for small~$m$.

\begin{table}[h]
\centering
\begin{tabular}{r r r r r r}
\toprule
$m$ & 0 blocks & 1 block & 2 blocks & 3+ blocks & Total \\
\midrule
1 & 1 & 1 & 0 & 0 & 2 \\
2 & 1 & 12 & 0 & 0 & 13 \\
3 & 1 & 108 & 5 & 0 & 114 \\
4 & 1 & 1{,}030 & 115 & 0 & 1{,}146 \\
5 & 1 & 10{,}785 & 1{,}764 & 28 & 12{,}578 \\
6 & 1 & 121{,}478 & 24{,}136 & 966 & 146{,}581 \\
7 & 1 & 1{,}442{,}392 & 320{,}844 & 20{,}877 & 1{,}784{,}114 \\
8 & 1 & 17{,}808{,}426 & 4{,}260{,}321 & 374{,}484 & 22{,}443{,}232 \\
9 & 1 & 226{,}558{,}935 & 57{,}015{,}530 & 6{,}147{,}306 & 289{,}721{,}772 \\
10 & 1 & 2{,}951{,}472{,}667 & 770{,}955{,}559 & 96{,}361{,}551 & 3{,}818{,}789{,}778 \\
\bottomrule
\end{tabular}
\caption{Sector decomposition of $|\CC([m]^2)|$ by number of contiguous blocks. The empty set (order-convex with no active rows) contributes 1 to the 0-block column.}
\label{tab:sectors}
\end{table}


%% ================================================================
\section{The Supermultiplicativity Theorem}\label{sec:supermult}
%% ================================================================

\begin{proof}[Proof of Theorem~\ref{thm:supermult}]
We construct an injection
\[
\CC([m]^2) \times \CC([n]^2) \;\hookrightarrow\; \CC([m+n]^2).
\]
Given convex $S_1 \subseteq [m]^2$ and $S_2 \subseteq [n]^2$, define embeddings
\begin{align*}
\varphi_1 &: [m]^2 \to [m+n]^2, \quad (i, j) \mapsto (i,\; j + n), \\
\varphi_2 &: [n]^2 \to [m+n]^2, \quad (i, j) \mapsto (i + m,\; j).
\end{align*}
Set $S := \varphi_1(S_1) \cup \varphi_2(S_2)$.

\textit{Both images are convex.}
The embeddings $\varphi_1, \varphi_2$ are order-preserving and order-reflecting (they preserve and reflect the component-wise partial order).
A convex subset mapped by an order-embedding remains convex, since any element between two image points must itself be an image point (by the bounds on row and column indices).

\textit{The images are anti-monotone.}
Every element of $\varphi_1(S_1)$ has column index $\geq n$ and row index $< m$.
Every element of $\varphi_2(S_2)$ has column index $< n$ and row index $\geq m$.
For $(i_1, j_1 + n) \in \varphi_1(S_1)$ and $(i_2 + m, j_2) \in \varphi_2(S_2)$:
\[
j_1 + n \geq n > j_2 \quad\text{(since } j_2 < n\text{)},
\]
so these elements are incomparable regardless of the row indices.

\textit{The union is convex.}
The union of two convex sets with no comparable pairs between them is convex: if $a \leq c \leq b$ with $a, b$ in the union, then $a$ and $b$ must lie in the \emph{same} image (otherwise they would be comparable across images), and convexity of that image gives $c$ in the same image.

\textit{The map is injective.}
Since $\varphi_1(S_1)$ occupies rows $[0, m)$ with columns $[n, m+n)$ and $\varphi_2(S_2)$ occupies rows $[m, m+n)$ with columns $[0, n)$, the two images are disjoint.
Thus $S_1$ and $S_2$ are uniquely recoverable from $S$.
\end{proof}

\begin{proof}[Proof of Corollary~\ref{cor:fekete}]
Theorem~\ref{thm:supermult} gives $a_{m+n} \geq a_m \cdot a_n$, so $\log a_m$ is superadditive.
By Fekete's lemma~\cite{fekete1923}, $\lim_{m \to \infty} (\log a_m) / m = \sup_m (\log a_m) / m$, hence $\lim a_m^{1/m} = \sup a_m^{1/m}$.
Finiteness: every order-convex $S \subseteq [m]^2$ is determined by the pair $(I(S), F(S))$ of its downward closure and upward closure, since $S = I(S) \cap F(S)$ by convexity.
The number of order ideals of $[m]^2$ is $\binom{2m}{m}$ (monotone lattice paths).
Hence $a_m \leq \binom{2m}{m}^2 \leq 16^m$, giving $\rho \leq 16$.
\end{proof}


%% ================================================================
\section{Transfer-Matrix Enumeration}\label{sec:transfer}
%% ================================================================

The Endpoint Classification (Theorem~\ref{thm:classification}) enables a dynamic programming algorithm that constructs convex subsets row by row.

\begin{definition}
A \emph{DP state} for $n$-column grids is a triple $(\mathbf{c}, L_{\min}, \mathrm{flag})$ where:
\begin{itemize}[nosep]
\item $\mathbf{c} = (c_0, \ldots, c_{n-1}) \in (\{0, \ldots, n-1\} \cup \{\infty\})^n$ is the \emph{constraint tuple}: $c_j = \min\{R_i : L_i \leq j,\; \text{row } i \text{ active in current block}\}$, or $\infty$ if no such row exists;
\item $L_{\min} \in \{0, \ldots, n-1\} \cup \{\infty\}$ is the minimum $L$ value across all active rows in all previous blocks;
\item $\mathrm{flag} \in \{\mathrm{true}, \mathrm{false}\}$ indicates whether the current row is in a contiguous block.
\end{itemize}
\end{definition}

\begin{proof}[Proof of Proposition~\ref{thm:transfer}]
By exhaustive enumeration (verified computationally for $n \leq 15$), the number of reachable states is exactly $\binom{n+2}{2} = (n+1)(n+2)/2$.
Each transition corresponds to choosing a row state (empty or an interval $[L, R]$), checking compatibility with the constraint tuple, and updating the state.
The computation of $|\CC([m] \times [n])|$ reduces to $m$ steps of matrix-vector multiplication with a $\binom{n+2}{2} \times \binom{n+2}{2}$ transition matrix.
\end{proof}

This yields the first 50 terms of the square-grid sequence; Table~\ref{tab:sequence} lists the first 10.

\begin{table}[h]
\centering
\begin{tabular}{r r r r r}
\toprule
$m$ & $|\CC([m]^2)|$ & $|\mathrm{Int}([m]^2)|$ & $a_m / a_{m-1}$ & $a_m / \mathrm{FC}(m)$ \\
\midrule
1 & 2 & 1 & --- & 1.000 \\
2 & 13 & 9 & 6.50 & 1.000 \\
3 & 114 & 36 & 8.77 & 1.000 \\
4 & 1{,}146 & 100 & 10.05 & 0.997 \\
5 & 12{,}578 & 225 & 10.98 & 0.999 \\
6 & 146{,}581 & 441 & 11.65 & 1.009 \\
7 & 1{,}784{,}114 & 784 & 12.17 & 1.024 \\
8 & 22{,}443{,}232 & 1{,}296 & 12.58 & 1.045 \\
9 & 289{,}721{,}772 & 2{,}025 & 12.91 & 1.071 \\
10 & 3{,}818{,}789{,}778 & 3{,}025 & 13.18 & 1.101 \\
\bottomrule
\end{tabular}
\caption{First 10 terms of $|\CC([m]^2)|$ with interval counts $|\mathrm{Int}([m]^2)| = (m(m+1)/2)^2$, consecutive ratios, and comparison with $\mathrm{FC}(m) = \frac{2}{6m+2}\binom{6m+2}{m}$.}
\label{tab:sequence}
\end{table}


%% ================================================================
\section{Comparison with Fuss--Catalan Numbers}\label{sec:catalan}
%% ================================================================

The generalized Fuss--Catalan numbers with parameters $(p, r) = (6, 2)$ are
\[
\mathrm{FC}(m) \;=\; \frac{2}{6m+2}\binom{6m+2}{m},
\]
giving the sequence $2, 13, 114, 1150, 12586, 145299, \ldots$ (OEIS A212071~\cite{oeis}).
The first three terms coincide exactly with $a_m$:
\[
a_1 = \mathrm{FC}(1) = 2, \quad a_2 = \mathrm{FC}(2) = 13, \quad a_3 = \mathrm{FC}(3) = 114.
\]
The sequences diverge at $m = 4$: $a_4 = 1146$ while $\mathrm{FC}(4) = 1150$.

With 50 terms, the asymptotic behavior is clear: the ratio $a_m / \mathrm{FC}(m)$ diverges \emph{exponentially}.
At $m = 30$, $a_m / \mathrm{FC}(m) \approx 2.72$; at $m = 50$, $a_m / \mathrm{FC}(m) \approx 8.55$.
The log of this ratio grows linearly in~$m$ with slope approximately $0.053$, indicating that $a_m$ has a strictly larger exponential growth rate than $\mathrm{FC}(m)$.

The Fuss--Catalan numbers have growth constant $6^6/5^5 = 46656/3125 \approx 14.9299$.
Richardson extrapolation of the consecutive ratios of $a_m$ yields estimates converging to the exact growth constant $\rho = 16$, which exceeds $6^6/5^5$ by approximately $7\%$.

\begin{remark}
The coincidence $a_m = \mathrm{FC}(m)$ for $m \leq 3$ is not accidental.
For small~$m$, the gapped sector is empty ($m \leq 2$) or negligible (5 subsets out of 114 at $m = 3$).
The contiguous sector, which dominates at small~$m$, appears to be counted by a Catalan-adjacent model.
The divergence at $m = 4$ is the point where the multi-block sector first contributes significantly.
\end{remark}


%% ================================================================
\section{The Partition Function Interpretation}\label{sec:partition}
%% ================================================================

The block decomposition admits a natural statistical-mechanical interpretation.

\begin{definition}
The \emph{block-weighted partition function} is
\[
Z_m(\alpha) \;:=\; \sum_{\substack{S \subseteq [m]^2 \\ S \text{ convex}}} e^{-\alpha \cdot k(S)},
\]
where $k(S)$ is the number of contiguous blocks of~$S$.
\end{definition}

At $\alpha = 0$, $Z_m(0) = |\CC([m]^2)|$ (unweighted counting).
As $\alpha \to \infty$, only the 0-block (empty set) and 1-block (contiguous) sectors survive.

\begin{definition}
The \emph{free energy per row} is $f(\alpha) := \lim_{m \to \infty} \frac{1}{m} \log Z_m(\alpha)$, when the limit exists.
\end{definition}

The supermultiplicativity theorem gives $f(0) = \log \rho \approx 2.77$.
The contiguous sector has its own growth rate $\rho_{\mathrm{contig}} \approx 12.8$, giving $f(\infty) \approx 2.55$.
The \emph{gap entropy} $f(0) - f(\infty) \approx 0.22$ measures the excess entropy contributed by multi-block configurations.

The block count $k(S)$ is not an arbitrary parameter: by Proposition~\ref{prop:blocks}, it equals the number of comparability-connected components of the active rows, a structural invariant forced by the classification theorem.

The Noetherian ratio relates to the partition function as
\[
\gamma([m]^2) \;=\; \frac{Z_m(0)}{|\mathrm{Int}([m]^2)|} \;=\; \frac{Z_m(0)}{(m(m+1)/2)^2}.
\]


%% ================================================================
\section{Formal Verification}\label{sec:lean}
%% ================================================================

Theorems~\ref{thm:classification}--\ref{thm:supermult}, Corollary~\ref{cor:fekete}, and Theorems~\ref{thm:rho}--\ref{thm:dimension} are formally verified using the Lean~4 proof assistant with Mathlib~\cite{mathlib}.
The verification covers the endpoint classification, the block decomposition, supermultiplicativity, the ideal/filter injection, the existence and exact value of the growth constant, and the dimension law for all $d \geq 2$.
Exact counts for $[m]^2$ with $m \leq 4$ are kernel-verified; values for $m \geq 5$ are computed by the transfer-matrix algorithm.
Source code is available at~\cite{leanrepo}.


%% ================================================================
\section{Conjectures}\label{sec:conjectures}
%% ================================================================

\begin{conjecture}\label{conj:gf}
The generating function $A(z) = \sum_{m \geq 1} |\CC([m]^2)|\, z^m$ is D-finite (satisfies a linear ODE with polynomial coefficients).
\end{conjecture}

\begin{theorem}\label{thm:rho}
The growth constant is $\rho = 16 = 2^4$ exactly.
The upper bound $\rho \leq 16$ follows from the ideal/filter injection (Corollary~\ref{cor:fekete}); the matching lower bound is proved by showing the injection is asymptotically tight~\cite{leanrepo}.
\end{theorem}

\begin{conjecture}\label{conj:subexp}
The subexponential factor satisfies
\[
|\CC([m]^2)| \;\sim\; C \cdot 16^m \cdot m^{-\alpha}
\]
for constants $C > 0$ and $\alpha > 0$.
\end{conjecture}

\begin{theorem}\label{thm:dimension}
For $d$-dimensional grids $[m]^d$ with $d \geq 2$, the sequence $|\CC([m]^d)|$ is supermultiplicative, the growth constant $\rho_d := \lim_{m \to \infty} |\CC([m]^d)|^{1/m}$ exists and is finite, and $\log|\CC([m]^d)| = \Theta(m^{d-1})$.
\end{theorem}

\begin{proof}
Supermultiplicativity for all $d \geq 2$ is proved by the same concatenation argument as Theorem~\ref{thm:supermult}.
The upper bound $|\CC([m]^d)| \leq |\mathrm{downsets}([m]^d)|^2$ and the tiling lower bound $|\CC([km]^d)| \geq |\CC([m]^d)|^{\Theta(k^{d-1})}$ together give $\log|\CC([m]^d)| = \Theta(m^{d-1})$.
Formally verified; see~\cite{leanrepo}.
\end{proof}

\begin{conjecture}\label{conj:bijection}
There exists a bijective encoding of order-convex subsets of $[m]^2$ by pairs of monotone boundary paths with anti-monotone gap constraints, giving a combinatorial interpretation of the counting sequence.
\end{conjecture}


%% ================================================================
\section{Discussion}\label{sec:discussion}
%% ================================================================

The central finding of this paper is strong numerical evidence that order-convex subsets of grid posets form a new exponential combinatorial class, with a growth constant distinct from all known Catalan/Raney families.
The structural explanation is the two-sector decomposition: contiguous blocks behave in a Catalan-like manner, while anti-monotone gaps introduce an independent combinatorial layer that changes the leading growth rate.

The coincidence $a_m = \mathrm{FC}(m)$ for $m \leq 3$ is explained by the fact that the gapped sector is empty or negligible for small~$m$.
The divergence at $m = 4$ marks the onset of the multi-block sector, which grows to comprise $23\%$ of the total by $m = 10$ and drives the exponential separation from the Fuss--Catalan curve.

The supermultiplicativity theorem, whose proof uses the same anti-monotone embedding mechanism that generates the new class, converts the growth constant from numerical evidence to a theorem-level statement.

Several directions remain:
\begin{enumerate}[nosep]
\item \emph{Determine the subexponential factor}: prove or disprove $|\CC([m]^2)| \sim C \cdot 16^m \cdot m^{-\alpha}$ and identify the exponent~$\alpha$.
\item \emph{Derive a recurrence or functional equation} from the transfer-matrix specification.
\item \emph{Extend to $[m] \times [n]$} and compute the bivariate generating function, which may reveal the true combinatorial mechanism more clearly than the square case.
\item \emph{Determine exact constants $c_d$ for $d \geq 3$}: the dimension law $\log|\CC([m]^d)| = \Theta(m^{d-1})$ is proved (Theorem~\ref{thm:dimension}), but the exact growth constants $\rho_d$ for $d \geq 3$ are open.
\end{enumerate}


\section*{Acknowledgments}

The formal verification builds on the Mathlib library~\cite{mathlib}.


\begin{thebibliography}{99}

\bibitem{birkhoff1967}
G.~Birkhoff,
\emph{Lattice Theory}, 3rd ed.,
AMS Colloquium Publications, vol.~25, 1967.

\bibitem{barnette2019}
B.~Barnette, W.~Nichols, and T.~Richmond,
The number of convex sets in a product of totally ordered sets,
\emph{Rocky Mountain J.\ Math.} \textbf{49} (2019), no.~2, 369--385.

\bibitem{bombelli1987}
L.~Bombelli, J.~Lee, D.~Meyer, and R.D.~Sorkin,
Spacetime as a causal set,
\emph{Phys.\ Rev.\ Lett.} \textbf{59} (1987), no.~5, 521--524.

\bibitem{difiore2026cag}
T.~DiFiore,
Causal-algebraic geometry: A Grothendieck-type framework for partial orders,
preprint, 2026.

\bibitem{dilworth1950}
R.P.~Dilworth,
A decomposition theorem for partially ordered sets,
\emph{Ann.\ of Math.\ (2)} \textbf{51} (1950), 161--166.

\bibitem{fekete1923}
M.~Fekete,
\"Uber die Verteilung der Wurzeln bei gewissen algebraischen Gleichungen mit ganzzahligen Koeffizienten,
\emph{Math.\ Z.} \textbf{17} (1923), 228--249.

\bibitem{gratzer2010}
G.~Gr\"atzer and R.W.~Quackenbush,
Positive universal classes in locally finite varieties,
\emph{Algebra Universalis} \textbf{64} (2010), 1--13.

\bibitem{mathlib}
The Mathlib Community,
\emph{Mathlib: The Lean Mathematical Library},
\url{https://github.com/leanprover-community/mathlib4}, 2024.

\bibitem{oeis}
OEIS Foundation Inc.,
\emph{The On-Line Encyclopedia of Integer Sequences},
\url{https://oeis.org}, 2024.

\bibitem{semenova2002}
M.V.~Semenova and F.~Wehrung,
Sublattices of lattices of order-convex sets, II: Posets of finite length,
\emph{Internat.\ J.\ Algebra Comput.} \textbf{13} (2003), no.~5, 543--564.

\bibitem{leanrepo}
T.~DiFiore,
\emph{Causal-Algebraic Geometry in Lean~4},
\url{https://github.com/tomdif/causal-algebraic-geometry-lean}, 2026.

\end{thebibliography}

\end{document}
